% Source: https://artofproblemsolving.com/community/c3h2435726_the_shoelace_theorem
% Stable identity: aops-topic-2435726-q01
\textbf{Step 1: a triangle with a vertex at the origin.} Let $O$ be the origin and let $P=(x_P,y_P)$, $Q=(x_Q,y_Q)$ be such that $O,P,Q$ are in counterclockwise order, i.e. $P=r(\cos\alpha,\sin\alpha)$ and $Q=s(\cos\beta,\sin\beta)$ with $0<\gamma=\beta-\alpha<180^\circ$. Then
\[[OPQ]=\dfrac12rs\sin\gamma=\dfrac12rs\left(\sin\beta\cos\alpha-\cos\beta\sin\alpha\right)=\dfrac12\left(x_Py_Q-x_Qy_P\right).\]
\textbf{Step 2: the polygon, with the origin inside.} Suppose first that the origin $O$ is an interior point of the polygon. The segments $OP_1,\dots,OP_n$ cut the convex polygon into the triangles $OP_iP_{i+1}$, each with its vertices in counterclockwise order. Adding Step 1 over these triangles gives
\[\text{area}=\sum_{i=1}^{n}[OP_iP_{i+1}]=\dfrac12\sum_{i=1}^{n}\left(x_iy_{i+1}-x_{i+1}y_i\right).\]

\begin{asy}
usepackage("amsmath");
size(5cm);
defaultpen(fontsize(10pt));
pair[] P={(0,0),(4,-0.5),(5.5,2),(3.5,4.5),(0.5,3.5)};
pair O=(2.6,1.8);
fill(O--P[1]--P[2]--cycle,rgb(1,0.88,0.7));
draw(P[0]--P[1]--P[2]--P[3]--P[4]--cycle,linewidth(1));
for(int i=0;i<5;++i) draw(O--P[i],dashed);
dot(O); label("$O$",O,dir(105));
string[] nm={"$P_1$","$P_2$","$P_3$","$P_4$","$P_5$"};
pair[] al={SW,SE,E,N,NW};
for(int i=0;i<5;++i){dot(P[i]); label(nm[i],P[i],al[i]);}
real a1=degrees(P[1]-O)-360, a2=degrees(P[2]-O);
draw(arc(O,0.6,a1,a2),red); label("$\gamma$",O+0.8*dir((a1+a2)/2),red);
\end{asy}

\textbf{Step 3: any position of the axes.} Moving the origin to an interior point means replacing every $(x_i,y_i)$ by $(x_i-p,\,y_i-q)$. The sum does not change:
\[\sum\left[(x_i-p)(y_{i+1}-q)-(x_{i+1}-p)(y_i-q)\right]=\sum\left(x_iy_{i+1}-x_{i+1}y_i\right)-q\sum(x_i-x_{i+1})-p\sum(y_{i+1}-y_i),\]
and the last two sums are $0$ because they telescope around the closed polygon. The area does not change either, so the formula holds for every choice of axes. $\blacksquare$

(The same formula, with an absolute value, holds for every simple polygon, convex or not; this needs signed areas of the triangles $OP_iP_{i+1}$, which can overlap and cancel.)
